\documentclass[12pt]{article}

\usepackage[portuguese]{babel}
\usepackage{float}
\usepackage{listings}
\usepackage{indentfirst}
\usepackage{amsfonts, amssymb, amsmath}
\usepackage{graphicx}
\usepackage{enumitem}
\usepackage[colorlinks=true, allcolors=blue]{hyperref}
\usepackage[letterpaper,
            top=1.5cm,
            bottom=1.5cm,
            left=2.75cm,
            right=2.75cm,
            marginparwidth=1.75cm]{geometry}

\title{Programação Orientada a Objetos}
\author{Prof. João Choma Neto}
\date{\today}

\begin{document}
\maketitle
GitHub do Professor: \url{https://github.com/JoaoChoma/}

\tableofcontents
\addcontentsline{toc}{section}{Apresentação da Matéria}
\section*{Apresentação da Matéria}
    \subsection{O que vamos estudar?}
    \begin{itemize}
        \item Programação orientada a objetos: abstração, encapsulamento, 
        herança e polimorfismo
        \item Implementação de aplicações em linguagem orientada a objetos
        \item Interface grafica do usuário e a programação orientada a eventos
        \item Presistência de objetos e integração com banco de dados
    \end{itemize}
    \subsection{Por que esta disciplina é importante?}
    \begin{itemize}
        \item Organização do código -- A orientação a objetos ajuda a 
        estruturar melhor responsabilidades, reduzir acoplamento e favorecer 
        manutenção.
        \item Modelagem mais próxima do problema -- Classes e objetos permitem 
        representar entidades, comportamentos e relacionamentos de forma mais 
        natural.
        \item Base do desenvolvimento moderno -- Muitas aplicações reais usam 
        POO em conjunto com interface, eventos, bancos de dados, frameworks.
        \item Aprendizagem progressiva -- O semestre evolui da base conceitual 
        para uma aplicação semestral mais completa.
        \item Integração de camadas -- A disciplina não para nas classes: ela 
        chega à interface grafica e à persistência de dados. 
    \end{itemize}

    \subsection{Fluxo do Semestre}
    \begin{enumerate}
        \item Fundamentos -- Classes, objetos, encapsulamento, heranla e 
        polimorfismo.
        \item Implementação -- Variaveis, estruturas, coleções, exceções, 
        classes abstratas e interfaces.
        \item Interface Gráfica -- MVC, componentes, layouts, e tratamento de 
        eventos.
        \item Persistência -- Acesso a banco, mapeamento objeto-relacional, 
        POJO e DAO.
        \item Projetos semestral -- Aplicação evoluída ao longo do semestre.
    \end{enumerate}

\section{Conceitos}
\begin{itemize}
    \item Objetos
    \begin{itemize}
        \item Qualquer elemento concreto ou abstrato que existe no mundo real, 
        que se pode individualizar por possuir comportamentos e caracteristicas 
        proprios.
        \item Qualquer coisa do mundo real que pode ser representada no mundo 
        virtual.
        \item Objetos abstratos
        \begin{itemize}
            \item Coisas que não são visiveis, uma ideia, como contas 
            bancárias e transações financeiras.
        \end{itemize}
        \item Objetos concretos
        \begin{itemize}
            \item Coisas que são palpaveis, como um carro ou uma pessoa.
        \end{itemize}
    \end{itemize}
    \item Classe
    \begin{itemize}
        \item Representa a abstração de um conjunto de objetos do mundo real 
        que possui comportamentos e caracteristicas comuns.
        \item Atributos: Caracteristicas particulares que os objetos de uma 
        classe possuem, assumindo valores diferentes para cada objeto.
        \item Métodos/Operações: Ordens que fazem o objeto executar ações. As 
        operações são tudo o que os objetos de uma calsse fazem. O método é a 
        lógica interna de uma operação. 
    \end{itemize}
    \item Instância
    \begin{itemize}
        \item É cada umas das ocorrências de um objeto formado a partir da 
        sua classe.
    \end{itemize}
    \item Atributos
    \begin{itemize}
        \item É uma caracteristica particular que os objetos de uma classe 
        possuem, assumindo valores diferentes para cada objeto.
    \end{itemize}
    \item Operações
    \begin{itemize}
        \item É uma ordem que faz um objeto executar uma ação. As operações são 
        tudo o que os objetos de uma classe fazem e nada além do que esses objetos 
        podem fazer.
    \end{itemize}
    \item Métodos
    \begin{itemize}
        \item É o algoritmo (conjunto de passos) que um objeto executa quando 
        reage a uma operação. O método é a lógica interna de uma operação.
        \item Dentro da programação orientada a objetos nós não acessamos as 
        variáveis diretamente, acessamos ela por meio de métodos.
    \end{itemize}
    
\end{itemize}

\end{document}